% manuscript.tex — Adhesion-controlled rigidity transition in a multi-cell phase field model
% Target: Physical Review E  (revtex4-2, two-column)
%
% Build:  pdflatex manuscript && pdflatex manuscript
%
\documentclass[aps,pre,twocolumn,superscriptaddress,longbibliography,floatfix]{revtex4-2}

% ── packages ──────────────────────────────────────────────────────────────────
\usepackage{amsmath,amssymb}
\usepackage{graphicx}
\usepackage{bm}              % bold math
\usepackage{xcolor}
\usepackage{hyperref}
\usepackage{booktabs}        % nice tables
\usepackage{lineno}           % line numbers for review
\linenumbers

% ── convenience macros ────────────────────────────────────────────────────────
\newcommand{\dd}{\mathrm{d}}               % upright differential
\newcommand{\vA}{v_{\!A}}                  % self-propulsion speed
\newcommand{\peff}{p_{\mathrm{eff}}}       % effective shape index
\newcommand{\Jt}{\tilde{J}}               % dimensionless adhesion J/(2γ)
\newcommand{\phii}{\phi_i}                 % cell field shorthand
\newcommand{\phij}{\phi_j}                 % neighbor field shorthand
\newcommand{\TODO}[1]{{\color{red}\textbf{[TODO: #1]}}}

% ══════════════════════════════════════════════════════════════════════════════
\begin{document}

\title{Adhesion-controlled rigidity transition in a multi-cell phase field model}

\author{Steven~A.~Silber}
\affiliation{Department of Physics and Astronomy, Western University, 1151 Richmond Street, London, Ontario, Canada N6A 3K7}

\author{Mikko~Karttunen}
\email{mikko.karttunen1@uef.fi}
\affiliation{Department of Physics and Astronomy, Western University, 1151 Richmond Street, London, Ontario, Canada N6A 3K7}
\affiliation{ELLIS Institute Finland, Maarintie 8, 02150 Espoo, Finland}
\affiliation{Department of Technical Physics, University of Eastern Finland, P.O.\ Box 1627, FI-70211 Kuopio, Finland}
\affiliation{Department of Chemistry, Western University, 1151 Richmond Street, London, Ontario, Canada N6A 5B7}

\date{\today}

% ── abstract ──────────────────────────────────────────────────────────────────
\begin{abstract}
The vertex model predicts a rigidity transition at a critical target shape index
$p_0^*\!\approx 3.81$, controlled by the balance between cortical contractility
and cell-cell adhesion.  Whether an analogous transition exists in continuum
phase field descriptions, where cell boundaries are diffuse and topology is
unconstrained, remains an open question.  In this work we add gradient-coupling
adhesion,
$F_{\mathrm{adh}} = J\sum_{i<j}\!\int\!\nabla\phii\cdot\nabla\phij\,\dd A$,
to the multi-cell phase field model of Palmieri
\textit{et al.}~\cite{Palmieri2015} and derive the
stability bound~$J < 2\gamma$.  The dimensionless parameter
$\Jt = J/(2\gamma)$ measures the fraction of surface energy removed at shared
interfaces.  Two-cell simulations confirm mass conservation, shape
preservation, and the predicted critical point.  Adhesion quench
experiments, in which adhesion is instantaneously enabled from an equilibrated
zero-adhesion state, show that cells undergo smooth, sub-cell interface
relaxation whose magnitude grows monotonically with~$\Jt$, but do not
undergo neighbor exchanges at zero motility.  Because the phase field
dynamics are continuous (unlike the discrete topology changes of the
vertex model), the quench traces the full relaxation trajectory from the
repulsive equilibrium to the adhesive one.  Motility probe experiments at
small~$\vA$ locate the adhesion-controlled jammed-to-fluid boundary.
The dimensionless parameter $\Jt$ thus plays the same role as the vertex
model target shape index $p_0$, enabling quantitative comparison
between the two frameworks.
\end{abstract}

\maketitle

% ══════════════════════════════════════════════════════════════════════════════
%  I.  INTRODUCTION
% ══════════════════════════════════════════════════════════════════════════════
\section{Introduction}
\label{sec:intro}

Epithelial monolayers can switch between a solid-like, mechanically rigid
state and a fluid-like, collectively migrating one.  This jamming transition
has been observed in wound-edge MDCK
monolayers~\cite{Angelini2011,Garcia2015}, in bronchial epithelial cells
from asthmatic donors~\cite{Park2015}, along the vertebrate body
axis~\cite{Mongera2018}, and in a variety of other epithelial
systems~\cite{Atia2018,Malinverno2017}.  The experimental evidence
consistently points to cell shape as the relevant order parameter: unjammed
monolayers contain elongated cells with a shape index $p = P/\sqrt{A}$
exceeding $\approx 3.81$, while jammed tissues consist of compact
cells~\cite{Park2015,Atia2018}.

The theoretical picture has been built largely on the vertex model, in
which each cell is a polygon with energy
$E = \sum_i[K_A(A_i - A_0)^2 + K_P(P_i - P_0)^2]$, and the target shape
index $p_0 = P_0/\sqrt{A_0}$ serves as the control
parameter~\cite{NagaiHonda2001,Farhadifar2007,Staple2010,Fletcher2014}.
Bi \textit{et al.}~\cite{Bi2015} showed that a rigidity transition occurs at
$p_0^* \approx 3.81$, and with added self-propulsion the transition extends
to a full $(p_0,\,v_0)$ phase diagram~\cite{Bi2016}.  Subsequent work
introduced active polarity~\cite{Barton2017}, multicellular
rosettes~\cite{YanBi2019}, 3D shape criteria~\cite{MerkelManning2018},
and cell turnover~\cite{Czajkowski2019}.  In the vertex model, adhesion
enters indirectly: cadherin-mediated adhesion reduces the line tension at
cell-cell junctions, which raises $p_0$ and pushes the tissue toward the
fluid phase.

Phase field models offer a complementary description in which each cell~$i$
is represented by a continuous scalar field $\phii(\bm{r},t)$, unity inside
the cell and zero outside, with a diffuse interface of width~$\lambda$.
Rearrangements occur through smooth field evolution rather than discrete
topology changes, and cell shapes are unconstrained by a polygonal
representation~\cite{Camley2017,AlertTrepat2020}.  Several multi-cell
formulations exist~\cite{Nonomura2012,Palmieri2015,Lober2015,Najem2016},
with various treatments of cell-cell adhesion
(Sec.~\ref{sec:model:literature}), but none has systematically varied
adhesion strength across a rigidity transition.

We add gradient-coupling adhesion
%
\begin{equation}
  F_{\mathrm{adh}} = J\sum_{i<j}\int \nabla\phii\cdot\nabla\phij\;\dd A,
  \label{eq:adhesion-intro}
\end{equation}
%
to the Palmieri \textit{et al.}~\cite{Palmieri2015} model, which uses
quartic repulsion but no explicit adhesion.  The dimensionless
control parameter $\Jt = J/(2\gamma)$ measures the fraction of surface
energy removed at shared interfaces, providing the phase field
counterpart of the vertex model target shape index~$p_0$.

We probe the transition using three protocols.
(i)~An \emph{adhesion quench} (Sec.~\ref{sec:quench}) at zero motility,
in which adhesion is switched on instantaneously from an equilibrated
repulsive state.  This protocol has no vertex model analog, since at
$T = 0$ the vertex model is frozen regardless of
$p_0$~\cite{Bi2015}; the phase field, by contrast, relaxes via gradient
descent, exposing the continuous trajectory of rearrangement events.
(ii)~A \emph{motility probe} (Sec.~\ref{sec:motility}), in which we sweep
$\Jt$ at small fixed $\vA$ to locate the jammed-to-fluid boundary.
(iii)~A full \emph{$(\Jt,\,\vA)$ phase diagram}
(Sec.~\ref{sec:diagram}), compared to the $(p_0, v_0)$ diagram of
Bi \textit{et al.}~\cite{Bi2016}.

The model and equations of motion are defined in Sec.~\ref{sec:model}
and simulation methods in Sec.~\ref{sec:methods}.


% ══════════════════════════════════════════════════════════════════════════════
%  II.  MODEL
% ══════════════════════════════════════════════════════════════════════════════
\section{Model}
\label{sec:model}

% --- II.A  Phase field representation and free energy ---
\subsection{Free energy}
\label{sec:model:energy}

We build on the multi-cell phase field model of Palmieri
\textit{et al.}~\cite{Palmieri2015}, in which $N$ cells occupy a
two-dimensional square domain of side~$L$ with periodic boundary conditions.
Each cell $i = 1,\dots,N$ is described by a scalar order parameter
$\phii(\bm{r},t) \in [0,1]$, where $\phii \approx 1$ inside the cell and
$\phii \approx 0$ outside.  Following Palmieri
\textit{et al.}~\cite{Palmieri2015}, the single-cell free energy contains
a gradient term, a double-well potential, and an area constraint;
cell-cell interactions add a quartic repulsion.  We extend this with
a gradient-coupling adhesion term, giving the total free energy
%
\begin{equation}
  F = F_{\mathrm{int}} + F_{\mathrm{vol}} + F_{\mathrm{rep}} + F_{\mathrm{adh}}.
  \label{eq:free-energy}
\end{equation}

The single-cell free energy contains a gradient and double-well
term~\cite{Palmieri2015}:
%
\begin{equation}
  F_{\mathrm{int}} = \sum_i \int \biggl[
    \gamma \, |\nabla\phii|^2
    + \frac{30\gamma}{\lambda^2}\,\phii^2(1-\phii)^2
  \biggr] \dd A,
  \label{eq:F-int}
\end{equation}
%
where $\gamma$ is the gradient energy coefficient (interfacial energy
scale) and $\lambda$ is the interface width.  The double-well potential
$V(\phi) = (30\gamma/\lambda^2)\,\phi^2(1-\phi)^2$ has minima at
$\phi = 0$ and $\phi = 1$ and sets the equilibrium interface profile to
a $\tanh$ shape of width~$\lambda$.  A soft volume constraint penalizes
deviations from the target area~$A_0 = \pi R^2$:
%
\begin{equation}
  F_{\mathrm{vol}} = \mu \sum_i
    \biggl(\int \phii^2 \, \dd A - A_0\biggr)^{\!2}.
  \label{eq:F-vol}
\end{equation}

Cell-cell interactions comprise steric repulsion and adhesion:
%
\begin{equation}
  F_{\mathrm{rep}} = \kappa \sum_{i<j} \int \phii^2\,\phij^2 \; \dd A,
  \label{eq:F-rep}
\end{equation}
%
\begin{equation}
  F_{\mathrm{adh}} = J \sum_{i<j} \int
    \nabla\phii \cdot \nabla\phij \; \dd A.
  \label{eq:F-adh}
\end{equation}

The quartic repulsion~$F_{\mathrm{rep}}$ prevents cell
overlap~\cite{Nonomura2012,Palmieri2015}.  Because $\phii^2\phij^2$ grows
as the fourth power of the fields, the cost of overlap rises steeply once
two cells' interiors begin to coincide.

The gradient-coupling adhesion~$F_{\mathrm{adh}}$ is the term added
in this work.  At a shared interface, the gradients of neighboring
cells are anti-parallel ($\nabla\phii \approx -\nabla\phij$), making
the integrand negative and reducing the energy.  Away from interfaces
both gradients vanish, so the coupling is surface-localized.

% --- II.B  Adhesion in the literature ---
\subsection{Context: adhesion in multi-cell phase field models}
\label{sec:model:literature}

Three distinct adhesion mechanisms have appeared in the multi-cell phase
field literature.

Nonomura~\cite{Nonomura2012} designed the first multi-cell phase field
model around the smooth step function $h(\phi) = \phi^2(3 - 2\phi)$, whose
gradient $\nabla h(\phii)$ is localized at the diffuse interface and
serves as a mathematical representation of the cell cortex.
Adhesion is a gradient-overlap coupling
$-\gamma_N\sum_{i<j}\!\int\!\nabla h(\phii)\!\cdot\!\nabla h(\phij)\,\dd A$,
supplemented by a regularization term
$c \sum_i \int |\nabla h(\phii)|^2 \, \dd A$ with $c > \gamma_N$.
The regularization adds gradient stiffness that stabilizes the interface
against adhesion-induced dissolution~\cite{Nonomura2012}, at the cost of a
second free parameter.  Using this model, Nonomura demonstrated
differential-adhesion-driven cell sorting in binary mixtures.

L\"ober, Ziebert, and Aranson~\cite{Lober2015} extended their single-cell
actin-substrate motility model~\cite{Ziebert2012} to many cells.  Because
the underlying model already contains non-variational terms (actin
advection, substrate friction, motor contractility), interactions are added
at the equation-of-motion level rather than through a free energy.  Their
adhesion term
$-\kappa_L\,\nabla\rho_i \!\cdot\! \hat{f}(\nabla\rho_j)$
advects cell~$i$'s boundary along the outward normal of cell~$j$,
modeling the mechanical attraction of adjacent cell
membranes~\cite{Lober2015}.  They used this to study
collision outcomes and collective migration of deformable cells.
This formulation cannot, by construction, be derived from a free energy
functional.

Najem and Grant~\cite{Najem2016} began from an explicit continuum
mechanics picture: an exponentially decaying adhesion potential
$W(x) = w\exp(-d^2/\delta^2)$ between cell surfaces.  They translated
this into phase field language via an auxiliary range field~$C_i$ that
defines each cell's interaction neighborhood independently of the
interface width.  The resulting adhesion energy
$w\,C_i\,\phii^2(1 - \phii)\sum_{j\neq i}(1 - \phij)$
has a range set by~$C_i$ rather than by the interface width, allowing
them to measure tissue surface tension as a function of the
adhesion-to-cortical-tension ratio.

The remaining multi-cell phase field
studies~\cite{Palmieri2015,Loewe2020,Wenzel2021,Graham2024,Saito2024} omit
adhesion entirely.  Palmieri \textit{et al.}~\cite{Palmieri2015} use
quartic repulsion only; Loewe \textit{et al.}~\cite{Loewe2020} study
the solid-liquid transition via deformability and motility;
Wenzel and Voigt~\cite{Wenzel2021} compare activity models;
Graham \textit{et al.}~\cite{Graham2024} demonstrate activity-driven
cell sorting; and Saito and Ishihara~\cite{Saito2024} show that
deformability alone can produce a fluid-to-fluid transition.

% --- II.C  Choice of the gradient-coupling form ---
\subsection{Choice of the gradient-coupling form}
\label{sec:model:choice}

Our formulation~\eqref{eq:F-adh} is the simplest member of the
Nonomura~\cite{Nonomura2012} family of gradient-overlap couplings.
Nonomura's original form uses the smooth-step function
$h(\phi) = \phi^2(3-2\phi)$ and requires a regularization term
$c\sum_i\!\int\!|\nabla h(\phii)|^2\dd A$ ($c > \gamma_N$)
to prevent adhesion-induced interface dissolution.  The
regularization widens the stable adhesion range but introduces a
second free parameter whose value must be chosen
empirically~\cite{Nonomura2012}.

We avoid the regularization entirely by coupling
$\nabla\phii$ directly rather than $\nabla h(\phii)$.
Physically, both forms are surface-localized:
both $\nabla\phii$ and $\nabla h(\phii)$ vanish in the cell
interior and exterior where $\phii$ is constant.  The quantitative
difference is a profile-dependent weighting factor
$h'(\phi) = 6\phi(1-\phi)$ that enters Nonomura's integrand
but not ours; the sharp-interface limit
(Sec.~\ref{sec:model:sharp}) shows that both forms produce an
adhesion energy proportional to the shared contact length.
The reason our simpler coupling does not require regularization
is that the stability bound $J < 2\gamma$
(Sec.~\ref{sec:model:stability}) is more restrictive than
Nonomura's regularized bound ($\gamma_N \lesssim 0.017$ vs.\
$2D_0 = 0.002$), confining $J$ to a range where the
interface remains stable without additional gradient stiffness.
The trade-off is a narrower stable adhesion range in
exchange for a single-parameter model with an analytically
derived stability criterion.

The variational derivative of the adhesion energy with respect to
cell~$i$ is
%
\begin{equation}
  \frac{\delta F_{\mathrm{adh}}}{\delta \phii}
    = -J \sum_{j \neq i} \nabla^2 \phij,
  \label{eq:adhesion-var}
\end{equation}
%
a Laplacian coupling: each cell feels the curvature of its neighbors'
fields.  Since $\nabla^2\phij$ is localized at cell~$j$'s interface (it
vanishes in the flat interior and exterior), the adhesion force is
automatically surface-localized without requiring the smooth step
function~$h$ or an auxiliary range field.

Unlike the non-variational advection of L\"ober
\textit{et al.}~\cite{Lober2015}, the gradient-coupling form derives from
a free energy functional.  This ensures that the $\vA = 0$ dynamics are
purely relaxational (gradient descent on~$F$), which is
required for the quench experiments of
Sec.~\ref{sec:quench}.

% --- II.D  Stability bound ---
\subsection{Stability bound}
\label{sec:model:stability}

The gradient coupling has a critical point beyond which shared
interfaces become energetically unstable.  Consider two cells in contact,
with anti-parallel gradients at the shared boundary
($\nabla\phi_1 \approx -\nabla\phi_2$).  The total gradient-type energy
per unit contact length is
%
\begin{align}
  E_{\mathrm{shared}} &= \int
    \bigl[\gamma|\nabla\phi_1|^2 + \gamma|\nabla\phi_2|^2
    + J\,\nabla\phi_1 \cdot \nabla\phi_2\bigr] \dd x
    \nonumber\\
  &= (2\gamma - J)\!\int\! |\nabla\phi|^2 \dd x,
  \label{eq:stability}
\end{align}
%
where each cell contributes $\gamma|\nabla\phi|^2$ to the gradient
energy at the shared boundary.  Without adhesion ($J = 0$) the cost
of the shared interface equals that of two free surfaces,
$2\gamma \int |\nabla\phi|^2 \dd x$; adhesion reduces this by a
fraction $J/(2\gamma)$.  The dimensionless control parameter is
therefore
%
\begin{equation}
  \Jt = \frac{J}{2\gamma},
  \label{eq:Jtilde}
\end{equation}
%
which measures the fraction of surface energy removed at shared
interfaces.  At~$\Jt = 1$ ($J = 2\gamma$), the shared interface
has zero energy cost; above this critical point the system gains energy
by creating more interface, leading to runaway instability and cell
merger.  In practice the quartic repulsion~$\kappa \phii^2 \phij^2$
provides additional stabilization that shifts the effective merger point
slightly above $J = 2\gamma$, but the interface is already effectively
dissolved at the critical point.

For our parameters ($\gamma = 1$, following Palmieri
\textit{et al.}~\cite{Palmieri2015}), the stability bound is
$J < 2$, giving a physical range $\Jt \in [0, 1)$.  Values up
to $\Jt = 0.75$ ($J = 1.5$, corresponding to 75\%\ surface
energy reduction) remain well within the stable regime.

The same bound follows from a positive-definiteness argument.
Define $S = \phi_1 + \phi_2$ and use the identity
$\nabla\phi_1\!\cdot\!\nabla\phi_2 = \tfrac{1}{2}(|\nabla S|^2
- |\nabla\phi_1|^2 - |\nabla\phi_2|^2)$.  The total gradient-type
energy becomes
$(\gamma - J/2)|\nabla\phi_1|^2 + (\gamma - J/2)|\nabla\phi_2|^2
+ (J/2)|\nabla S|^2$,
which is positive-definite if and only if $\gamma > J/2$, i.e.,
$J < 2\gamma$.

\textit{Comparison with Nonomura.}---Nonomura's
model~\cite{Nonomura2012} uses gradient-energy coefficient
$D_0 = 0.001$, adhesion coupling $\gamma_N = 0.003$--$0.0065$, and
regularization $c = 0.01$.  Without regularization, the model would
require $\gamma_N < 2D_0 = 0.002$; the regularization raises
the effective stability bound to $\gamma_N \lesssim 0.017$ by
adding gradient stiffness proportional to $c\langle h'^2\rangle$
(where $\langle h'^2\rangle \approx 0.77$ is the profile-averaged
square of $h' = 6\phi(1-\phi)$), roughly an $8.7\times$ increase
over the unregularized bound.  Nonomura's strongest
two-dimensional adhesion achieves~${\sim}36\%$ surface energy reduction,
comparable to our $\Jt \approx 0.36$ ($J \approx 0.7$).  Our model
without regularization has a more restricted range but still supports
physically meaningful adhesion strengths up to 75\%\ surface reduction.

% --- II.E  Equations of motion ---
\subsection{Equations of motion}
\label{sec:model:eom}

The dynamics of each cell field follow from the model of Palmieri
\textit{et al.}~\cite{Palmieri2015}, which falls into the model-A
classification of Hohenberg and Halperin~\cite{HohenbergHalperin1977}
with an additional convective term.  The equation of motion is
%
\begin{equation}
  \frac{\partial \phii}{\partial t}
    + \bm{v}_i \cdot \nabla\phii
    = -M\,\frac{\delta F}{\delta \phii},
  \label{eq:eom}
\end{equation}
%
where $M = 1/2$ is the mobility and $\bm{v}_i$ decomposes into an
inactive part (determined by thermodynamic relaxation) and active
self-propulsion $\vA^{(i)}\hat{\bm{p}}_i$~\cite{Palmieri2015}.
The polarity $\hat{\bm{p}}_i$ reorients via a run-and-tumble
process with persistence time~$\tau$.  In the zero-motility limit
($\vA = 0$) the dynamics reduce to relaxation on~$F$.
The full variational derivative, including the adhesion
contribution $-J(\nabla^2 S - \nabla^2\phii)$ where
$S = \sum_k\phi_k$, is given in the Supplemental Material.

% --- II.F  Sharp-interface limit ---
\subsection{Sharp-interface limit}
\label{sec:model:sharp}

To connect the gradient coupling to vertex model adhesion, we evaluate
the adhesion integral for $\tanh$-profile
interfaces of width~$\lambda$ at a flat shared boundary of
length~$\ell_{ij}$.  Using
$\nabla\phi_1 = -(2\lambda)^{-1}\mathrm{sech}^2(x/\lambda)\,\hat{x}$
and $\nabla\phi_2 = +(2\lambda)^{-1}\mathrm{sech}^2(x/\lambda)\,\hat{x}$,
the normal integral evaluates to
$\int\! \nabla\phi_1\!\cdot\!\nabla\phi_2\,dx = -1/(3\lambda)$
(via $\int\!\mathrm{sech}^4 u\,du = 4/3$).  Summing over all shared
boundaries:
%
\begin{equation}
  F_{\mathrm{adh}} \sim -\frac{J}{3\lambda}\sum_{i<j}\ell_{ij},
  \label{eq:sharp-interface}
\end{equation}
%
which has the same form as the vertex model adhesion
$-\gamma_{\mathrm{vm}}\sum_{i<j}\ell_{ij}$ with
$\gamma_{\mathrm{vm}} = J/(3\lambda)$.  The adhesion energy is
proportional to the total shared contact length, as in the
vertex model.  The effective shape index
$\peff = P_i/\sqrt{A_i}$, measured from the $\phii = 0.5$ contour,
is the corresponding phase field observable.


% ══════════════════════════════════════════════════════════════════════════════
%  III.  SIMULATION METHODS
% ══════════════════════════════════════════════════════════════════════════════
\section{Simulation methods}
\label{sec:methods}

% --- III.A  Implementation ---
\subsection{GPU-accelerated implementation}
\label{sec:methods:gpu}

The model is implemented as a custom CUDA C++ code.  All $N$ phase fields
are stored on a single GPU and evolved simultaneously.  The per-cell fields
$\phii$ are represented on a uniform Cartesian grid of $L \times L$ points
with spacing $\Delta x = 1$.  At each time step a \emph{scatter} kernel
accumulates the sum fields
%
\begin{equation}
  S(\bm{r}) = \textstyle\sum_k \phi_k(\bm{r}), \qquad
  S_2(\bm{r}) = \textstyle\sum_k \phi_k^2(\bm{r})
  \label{eq:sum-fields}
\end{equation}
%
over all cells, and then a \emph{fused} kernel evaluates the
variational derivative~\eqref{eq:eom}, self-propulsion, and
forward-Euler update for every cell and grid point in a single pass.
The repulsive force on cell~$i$ is
$2\kappa\,\phii\bigl(S_2 - \phii^2\bigr)$; the adhesive force is
$-J\bigl(\nabla^2 S - \nabla^2\phii\bigr)$, computed via a
McLellan nine-point isotropic Laplacian stencil on~$S$ (matching the
stencil used for $\nabla^2\phii$).  When $J = 0$, the Laplacian of~$S$ is neither
computed nor applied, so the code incurs zero overhead for the purely
repulsive case.

All production runs were performed on NVIDIA H100~GPUs (Digital Research
Alliance of Canada).

% --- III.B  Parameters ---
\subsection{Simulation parameters}
\label{sec:methods:params}

Table~\ref{tab:params} summarizes the physical and numerical parameters,
which are held fixed across all adhesion experiments.

\begin{table}[b]
  \caption{%
    Fixed simulation parameters.  The time scale is set by $\tau = 10{,}000$
    (tumble time); the length scale by the cell radius $R = 49$.  The packing
    fraction $\phi = N \pi R^2 / L^2 = 0.89$ places the system in the
    confluent regime.
  }
  \label{tab:params}
  \begin{ruledtabular}
  \begin{tabular}{lll}
    Parameter & Symbol & Value \\ \hline
    Cell count       & $N$        & 288 (or 1152)     \\
    Domain side      & $L$        & 1562 (or 3124)    \\
    Cell radius      & $R$        & 49                \\
    Packing fraction & $\phi$     & 0.89              \\
    Interface width  & $\lambda$  & 7                 \\
    Gradient energy  & $\gamma$   & 1                 \\
    Repulsion        & $\kappa$   & 10                \\
    Volume penalty   & $\mu$      & 1                 \\
    Friction         & $\xi$      & 1500              \\
    Mobility         & $M$        & 0.5               \\
    Time step        & $\Delta t$ & 0.01              \\
    Tumble time      & $\tau$     & $10{,}000$        \\
    Boundary cond.   &            & periodic          \\
  \end{tabular}
  \end{ruledtabular}
\end{table}

The packing fraction $\phi = 0.89$ was chosen to match the near-confluent
regime in which the vertex model is defined.  At this density, cells tile
the domain with only small interstitial gaps, consistent with a confluent
monolayer.

% --- III.C  Equilibration ---
\subsection{Equilibration protocol}
\label{sec:methods:eq}

All adhesion experiments start from a well-equilibrated, purely repulsive
configuration.  One hundred independent realizations of 288 cells were
initialized at random positions on the $L = 1562$ domain and relaxed at
$\vA = 0$, $J = 0$ for $t = 80{,}000$ ($8\tau$), following the
equilibration protocol of Palmieri \textit{et al.}~\cite{Palmieri2015}.
At this point, all transient stress from the random initialization has
decayed and the system sits in a local energy minimum of the repulsive
free energy.  The resulting 100 checkpoints serve as independent starting
states for the adhesion experiments.

We verified convergence by confirming that control runs ($\vA = 0$,
$J = 0$) starting from these equilibrated states produce mean centroid
displacements below $0.02R$ over $20{,}000$ time units, negligible
compared to the cell radius~$R$.

% --- III.D  Adhesion quench ---
\subsection{Adhesion quench protocol}
\label{sec:methods:quench}

The adhesion quench probes the rigidity of the energy landscape by
instantaneously enabling adhesion in an equilibrated many-cell system.
Starting from an equilibrated $N = 288$ configuration
(Sec.~\ref{sec:methods:eq}), we sweep
$\Jt$ at adhesion strength $J \in \{0,\, 0.9375,\, 1.875,\, 2.8125\}$
(corresponding to
$\Jt \in \{0,\, 0.125,\, 0.250,\, 0.375\}$)
and evolve for $t = 20{,}000$ time units ($= 2\tau$).  With no
self-propulsion, the dynamics are purely overdamped relaxation on the
free energy landscape.  The total centroid displacement
$\sum_i |\Delta\bm{r}_i|$ is the primary diagnostic.

Only phase field models can perform this experiment.  In vertex models,
changing adhesion modifies~$p_0$, but the system can only explore
configuration space through discrete T1~events or explicit motility.
The overdamped phase field dynamics explore the landscape continuously,
revealing the distinction between rigid and fluid energy landscapes
without any activity.

% --- III.F  Observables ---
\subsection{Observables}
\label{sec:methods:obs}

\paragraph{Mean-squared displacement (MSD).}
The ensemble- and time-averaged MSD is
%
\begin{equation}
  \langle \Delta r^2(t) \rangle
    = \biggl\langle \frac{1}{N}\sum_{i=1}^{N}
      |\bm{r}_i(t_0 + t) - \bm{r}_i(t_0)|^2 \biggr\rangle_{t_0},
  \label{eq:msd}
\end{equation}
%
with displacements unwrapped across periodic boundaries.

\paragraph{Self-overlap function.}
The fraction of cells that have not displaced beyond a cage
radius~$a$ after time~$\Delta t$:
%
\begin{equation}
  Q(\Delta t) = \frac{1}{N}\sum_{i=1}^{N}
  \Theta\!\bigl(a - |\bm{r}_i(t_0 + \Delta t) - \bm{r}_i(t_0)|\bigr),
  \label{eq:Qt}
\end{equation}
%
where $a = 0.3\,d_{\mathrm{cell}} \approx 28$ grid units.  We fit
$Q(t)$ to a stretched exponential
$Q(t) = \exp[-(t/\tau_\alpha)^\beta]$, where $\beta < 1$ indicates a
broad distribution of relaxation times.

\paragraph{Non-Gaussian parameter.}
%
\begin{equation}
  \alpha_2(t)
    = \frac{\langle \Delta r^4(t)\rangle}
           {2\,\langle \Delta r^2(t)\rangle^2} - 1
  \label{eq:alpha2}
\end{equation}
%
measures dynamic heterogeneity; a peak in $\alpha_2$ signals a population
of fast-moving cells coexisting with caged cells.

\paragraph{Four-point susceptibility.}
%
\begin{equation}
  \chi_4(\Delta t) = N\bigl[\langle Q(\Delta t)^2 \rangle
  - \langle Q(\Delta t)\rangle^2\bigr],
  \label{eq:chi4}
\end{equation}
%
which measures the number of cells participating in cooperative
rearrangements.

\paragraph{Effective shape index.}
From VTK snapshots, we extract the $\phii = 0.5$ contour for each cell
using a marching-squares algorithm, then compute the perimeter $P_i$
and area $A_i$ of the contour:
%
\begin{equation}
  \peff = \frac{P_i}{\sqrt{A_i}}.
  \label{eq:peff}
\end{equation}

\paragraph{Energy decomposition.}
The free energy~\eqref{eq:free-energy} is decomposed into gradient,
bulk, repulsive, and adhesive contributions, each tracked as a function
of time.  In the quench protocol, the time series of each component
distinguishes interface relaxation (smooth gradient energy decay) from
topological rearrangements (stepwise drops in interaction energy).

% --- III.G  Protocol summary ---
\subsection{Experimental protocols}
\label{sec:methods:protocols}

Three experimental phases, summarized in Table~\ref{tab:protocols}, probe the
rigidity transition at increasing levels of detail.

\begin{table}[t]
  \caption{%
    Summary of the three experimental protocols.  All start from the same
    pool of 100 independently equilibrated checkpoints at $t = 80{,}000$.
  }
  \label{tab:protocols}
  \begin{ruledtabular}
  \begin{tabular}{lccc}
    & Phase 0     & Phase 1       & Phase 2           \\
    & (quench)    & (probe)       & (diagram)         \\ \hline
    $\Jt$
      & 0--0.75 (4+ pts)  & 0--0.75 (6 pts) & 0--0.75 (6 pts) \\
    $\vA$
      & 0               & 0.002          & 0.002--0.012    \\
    Duration
      & $2\tau$          & $5\tau$        & $5\tau$         \\
    Replicates
      & 1 (deterministic)& 3             & 3--10           \\
    Key observable
      & displacement,    & MSD,           & MSD,            \\
      & energy           & $\alpha_2$     & $\peff$         \\
  \end{tabular}
  \end{ruledtabular}
\end{table}


% ══════════════════════════════════════════════════════════════════════════════
%  IV.  RESULTS — Two-cell validation
% ══════════════════════════════════════════════════════════════════════════════
\section{Two-cell validation}
\label{sec:2cell}

To characterize the gradient-coupling adhesion in isolation, we
place two cells in a $400 \times 400$ domain ($R = 49$, $\vA = 0$)
at center-to-center distance $d_0 \approx 1.9R$ and evolve for
$t = 2\tau$ at each adhesion strength.
Figure~\ref{fig:twocell} shows equilibrium snapshots at four
representative values of~$\Jt$, and the contact angle~$\theta$
as a function of~$\Jt$.

\begin{figure}[t]
  \centering
  \includegraphics[width=0.24\columnwidth]{figures/fig1_snap_Jt0.12}%
  \includegraphics[width=0.24\columnwidth]{figures/fig1_snap_Jt0.25}%
  \includegraphics[width=0.24\columnwidth]{figures/fig1_snap_Jt0.50}%
  \includegraphics[width=0.24\columnwidth]{figures/fig1_snap_Jt0.75}\\[4pt]
  \includegraphics[width=0.95\columnwidth]{figures/fig1_contact_angle}
  \caption{Two-cell adhesion validation
    ($\gamma = 1$, $R = 49$, $\vA = 0$, $t = 2\tau$).
    Top: composite field $\sum_i\phi_i$ at equilibrium for
    $\Jt = 0.125$, $0.25$, $0.50$, and $0.75$, with
    per-cell $\phi_i = 0.5$ contours (white).
    Bottom: measured contact angle $\theta$ versus $\Jt$.
    The contact angle increases monotonically with adhesion
    strength, confirming that the gradient coupling produces
    the expected interfacial geometry.}
  \label{fig:twocell}
\end{figure}

The contact angle increases monotonically from
$\theta = 22.5^\circ$ at $\Jt = 0.125$ to $50.8^\circ$ at
$\Jt = 0.75$ (Fig.~\ref{fig:twocell}).  At low adhesion the
cells remain nearly circular with a narrow contact zone; at
$\Jt = 0.75$ the contact region flattens substantially and the
cells deform into a peanut-like shape with a well-defined
contact angle.


% ══════════════════════════════════════════════════════════════════════════════
%  V.  RESULTS — Adhesion quench
% ══════════════════════════════════════════════════════════════════════════════
\section{Adhesion quench at zero motility}
\label{sec:quench}

\subsection{Protocol}
\label{sec:quench:protocol}

Starting from a checkpoint equilibrated at $J = 0$, $\vA = 0$ for
$t = 80{,}000$ ($8\tau$), we instantaneously set $J > 0$ while keeping
$\vA = 0$ and evolve for an additional $2\tau$.  The dynamics are
purely relaxational: $\partial_t \phii = -M\,\delta F/\delta\phii$.
Because the $J = 0$ equilibrium is generally not a minimum of the
$J > 0$ energy landscape, the system flows downhill to a new minimum.

This ``$J$-quench'' has no vertex model analog.  In the vertex model at
$T = 0$, $v_0 = 0$, the system is frozen in a local minimum regardless
of~$p_0$; changing $p_0$ relabels the energy without inducing motion.
In the phase field, turning on adhesion creates an energetic driving
force and the overdamped dynamics trace out the continuous
relaxation path---including intermediate configurations inaccessible in
a discrete-topology model.

\subsection{Centroid displacement}
\label{sec:quench:displacement}

\begin{table}[b]
\caption{Phase~0 adhesion quench results: mean centroid displacement
  after $t = 20{,}000$ TU ($2\tau$) of overdamped relaxation from an
  equilibrated $N = 288$ configuration at $\vA = 0$.  All displacements
  are sub-cell, reflecting interface reshaping rather than neighbor
  exchanges.}
\label{tab:quench}
\begin{ruledtabular}
\begin{tabular}{ccccc}
$J/\kappa$ & $\Jt$ & $\langle\Delta r\rangle / R$ & RMS disp. & $v_{\mathrm{rms}}$ (final) \\
\midrule
0.00 & 0    & 0.025 & 1.2 & $4.4\times10^{-6}$ \\
0.05 & 0.25 & 0.058 & 2.8 & $6.5\times10^{-5}$ \\
0.10 & 0.50 & 0.086 & 4.2 & $3.7\times10^{-4}$ \\
0.15 & 0.75 & 0.132 & 6.5 & $2.2\times10^{-3}$ \\
\end{tabular}
\end{ruledtabular}
\end{table}

The quench results (Table~\ref{tab:quench}) show a smooth, monotonic
increase in displacement with~$\Jt$.  The $J = 0$ control produces
mean displacements of $0.025R$, negligible compared to the cell
radius, confirming that the starting state is fully relaxed.  At
$\Jt = 0.25$, cells displace by $0.058R$; at $\Jt = 0.75$, by
$0.132R$.  All displacements are sub-cell ($\ll R$), and no neighbor
exchanges (T1-like topology changes) occur at any adhesion strength.

The displacements reflect interface reshaping: contact angles adjust
and shared contact areas grow as adhesion lowers the effective surface
tension at cell-cell boundaries.  The smooth, monotonic trend is
consistent with the system relaxing to the nearest local energy
minimum while preserving the neighbor topology inherited from
the~$J = 0$ equilibrium.  At $\Jt = 0.75$, the final
$v_{\mathrm{rms}} = 2.2\times10^{-3}$ indicates the system has not
fully equilibrated within $2\tau$, suggesting the relaxation timescale
grows with~$\Jt$, consistent with proximity to a soft mode.

The absence of neighbor rearrangements is physically expected.  At the
confluent packing fraction $\phi = 0.89$, the saddle-point energy for
a topology change is dominated by the cell-squeezing cost ($\kappa$
repulsion and $\gamma$ gradient energy required to push one cell past
another), which adhesion does not reduce.  Adhesion lowers the energy
of configurations with more shared contact area (the destination) but
not the barrier.  At $\vA = 0$, gradient descent cannot overcome
these barriers regardless of~$\Jt$.

The quench thus measures the static equilibrium
properties at each adhesion strength.  The active driving needed to
overcome topological barriers is provided by motility, which is the
subject of Phase~1 (Sec.~\ref{sec:motility}).

\TODO{Fig: (a)~Centroid displacement $\langle\Delta r\rangle / R$ vs.\ $\Jt$ from
the adhesion quench (Table~\ref{tab:quench}).
(b)~Representative VTK snapshots before and after quench.
(d)~Energy time series during quench.}

\subsection{Energy landscape}
\label{sec:quench:energy}

\TODO{Energy time series for representative $\Jt$ values.  Expected:
smooth monotonic decay below $\Jt^*$; stepwise or multi-plateau decay
above $\Jt^*$ corresponding to discrete rearrangement events.}

\subsection{Neighbor rearrangements}
\label{sec:quench:topology}

\TODO{Count T1-like neighbor changes vs.\ $\Jt$.}


% ══════════════════════════════════════════════════════════════════════════════
%  VI.  RESULTS — Phase 1: Motility probe
% ══════════════════════════════════════════════════════════════════════════════
\section{Motility probe sweep}
\label{sec:motility}

At fixed $\vA = 0.002$, all tested adhesion strengths
($\Jt \leq 0.25$) remain jammed. The clean baseline
($\Jt = 0$) has $D_{\mathrm{eff}} = 6.4 \times 10^{-4}$
and $\tau_\alpha > 100{,}000$ (Q does not decay to $e^{-1}$
within $20\tau$), confirming the system is deep in the solid
phase. At $\Jt = 0.125$, $D_{\mathrm{eff}}$ decreases
slightly to $2.8 \times 10^{-4}$, indicating that weak
adhesion does not fluidize the system at this low motility.

\subsection{Mean-squared displacement}
\label{sec:motility:msd}

\TODO{MSD vs.\ time for each $\Jt$ at $\vA = 0.002$.}

\subsection{Dynamic heterogeneity}
\label{sec:motility:alpha2}

\TODO{$\alpha_2(t)$ for each $\Jt$.}

\subsection{Effective shape index}
\label{sec:motility:shape}

\TODO{Extract $\peff$ from trajectory $L_n$ column. Plot
$\langle\peff\rangle$ vs.\ $\Jt$; compare to
$p_0^* \approx 3.81$.}


% ══════════════════════════════════════════════════════════════════════════════
%  VII.  RESULTS — Phase 2: Phase diagram
% ══════════════════════════════════════════════════════════════════════════════
\section{Adhesion--motility phase diagram}
\label{sec:diagram}

Table~\ref{tab:phase2} summarizes the effective diffusion coefficient
$D_{\mathrm{eff}}$ and structural relaxation time $\tau_\alpha$
across the $(\Jt,\,\vA)$ parameter space.  Results for
$\Jt \geq 0.25$ are in progress at $\Delta t = 0.01$.

\begin{table}[t]
\caption{Phase~2 results: effective diffusion coefficient
  $D_{\mathrm{eff}}$ and relaxation time $\tau_\alpha$ (from
  stretched-exponential fit to $Q(t)$) for the adhesion--motility
  sweep.  $N = 288$, $\phi = 0.89$, production time $20\tau$ from
  equilibrated checkpoint.  Three independent replicates per point.
  $\beta = 2.0$ indicates Q did not decay sufficiently for a reliable
  stretched-exponential fit (jammed regime).}
\label{tab:phase2}
\begin{ruledtabular}
\begin{tabular}{ccccr}
$\Jt$ & $\vA$ & $D_{\mathrm{eff}}$ & $\tau_\alpha$ & $\beta$ \\
\midrule
\multicolumn{5}{c}{$\Jt = 0$ (no adhesion)} \\
0 & 0.002 & $4.8\times10^{-4}$ & $>10^5$ & 2.0 \\
0 & 0.004 & $8.8\times10^{-4}$ & $>10^5$ & 2.0 \\
0 & 0.006 & $1.5\times10^{-3}$ & $>10^5$ & 2.0 \\
0 & 0.008 & $6.3\times10^{-3}$ & $80{,}977$ & 1.24 \\
0 & 0.010 & $1.4\times10^{-2}$ & $47{,}761$ & 1.06 \\
0 & 0.012 & $2.3\times10^{-2}$ & $29{,}986$ & 0.88 \\
\midrule
\multicolumn{5}{c}{$\Jt = 0.125$ (12.5\% surface energy reduction)} \\
0.125 & 0.002 & $1.9\times10^{-4}$ & $>10^5$ & 2.0 \\
0.125 & 0.004 & $5.6\times10^{-4}$ & $>10^5$ & 2.0 \\
0.125 & 0.006 & $1.6\times10^{-3}$ & $99{,}433$ & 2.0 \\
0.125 & 0.008 & $7.1\times10^{-3}$ & $64{,}838$ & 1.23 \\
0.125 & 0.010 & $1.2\times10^{-2}$ & $36{,}197$ & 0.89 \\
0.125 & 0.012 & $3.4\times10^{-2}$ & $21{,}875$ & 0.89 \\
\end{tabular}
\end{ruledtabular}
\end{table}

The data show two trends. First, the clean-system ($\Jt = 0$)
transition is visible: $D_{\mathrm{eff}}$ increases by an order of
magnitude between $\vA = 0.006$ and $0.010$, and $\tau_\alpha$ drops
from $>10^5$ to $47{,}761$ ($\approx 5\tau$) with $\beta$ falling
below~1.
This locates the unjamming transition at $v_A^* \approx 0.009$ for
$\Jt = 0$, consistent with the Palmieri~\textit{et al.}~\cite{Palmieri2015}
calibration.

Second, adhesion shifts the transition.  At $\Jt = 0.125$,
$\tau_\alpha$ at $\vA = 0.008$ drops from $80{,}977$ to $64{,}838$
and $D_{\mathrm{eff}}$ increases from $6.3 \times 10^{-3}$ to
$7.1 \times 10^{-3}$.  At $\vA = 0.012$, $D_{\mathrm{eff}}$
increases by 50\% (from $2.3 \times 10^{-2}$ to
$3.4 \times 10^{-2}$) and $\tau_\alpha$ decreases from
$29{,}986$ to $21{,}875$.

\subsection{Phase boundary}
\label{sec:diagram:boundary}

\TODO{Heatmap of $D_{\mathrm{eff}}$ in the $(\Jt,\,\vA)$ plane
once $\Jt \geq 0.25$ data with McLellan stencil is available.
Compare to the $(p_0,\,v_0)$ boundary of Bi \textit{et al.}~\cite{Bi2016}.}

\subsection{Shape index mapping}
\label{sec:diagram:shape}

\TODO{Overlay contours of constant $\peff$ on the phase diagram.}

\subsection{System-size effects}
\label{sec:diagram:size}

\TODO{Compare 288-cell and 1152-cell results.}


% ══════════════════════════════════════════════════════════════════════════════
%  VIII.  DISCUSSION
% ══════════════════════════════════════════════════════════════════════════════
\section{Discussion}
\label{sec:discussion}

We have added gradient-coupling adhesion to the Palmieri
\textit{et al.}~\cite{Palmieri2015} multi-cell phase field model.
The dimensionless parameter $\Jt = J/(2\gamma)$ controls the
fraction of surface energy removed at shared interfaces, and the
sharp-interface limit (Sec.~\ref{sec:model:sharp}) confirms that
$\Jt$ is the phase field counterpart of the vertex model target
shape index~$p_0$.

The two-cell validation (Fig.~\ref{fig:twocell}) and the adhesion
quench (Table~\ref{tab:quench}) establish that the coupling is
well-behaved within the stability bound $J < 2\gamma$: mass is
conserved, cell shapes are preserved, and the quench produces
smooth interface relaxation without neighbor exchanges at
$\vA = 0$.  The quench protocol has no vertex model analog, since
at zero motility the vertex model is frozen regardless of $p_0$;
the phase field dynamics trace the full relaxation trajectory
from the repulsive to the adhesive equilibrium.

The motility sweep (Table~\ref{tab:phase2}) shows that adhesion
shifts the unjamming transition.  At $\Jt = 0.125$,
$\tau_\alpha$ at $\vA = 0.008$ drops by 20\% relative to $\Jt = 0$,
and at $\vA = 0.012$ $D_{\mathrm{eff}}$ increases by 50\%.
The clean-system transition at $v_A^* \approx 0.009$ is consistent
with the Palmieri \textit{et al.}~\cite{Palmieri2015} calibration.

Several questions remain open.  Extracting the effective shape
index $\peff$ from $\phii = 0.5$ contours will test whether the
transition at $\Jt^*$ coincides with
$\langle\peff\rangle \approx 3.81$.  The diffuse interface
(width $\lambda = 7$ vs.\ cell radius $R = 49$) may introduce
corrections relative to the sharp-polygon limit.  Comparing
$N = 288$ and $N = 1152$ will quantify finite-size effects, and
extension to three dimensions is direct since the gradient coupling
and the bound $J < 2\gamma$ hold in any spatial dimension.


% ══════════════════════════════════════════════════════════════════════════════
%  ACKNOWLEDGMENTS
% ══════════════════════════════════════════════════════════════════════════════
\begin{acknowledgments}
Computational resources were provided by the Digital Research Alliance
of Canada.
M.K.\ thanks the Natural Sciences and Engineering Research Council of
Canada (NSERC), the Canada Research Chairs Program, and Foundation PS
for financial support.
S.A.S.\ thanks NSERC for financial support through the Postgraduate
Scholarships (PGS~D) program.
\end{acknowledgments}

% ══════════════════════════════════════════════════════════════════════════════
%  BIBLIOGRAPHY
% ══════════════════════════════════════════════════════════════════════════════
\begin{thebibliography}{99}

\bibitem{Angelini2011}
T.~E. Angelini \textit{et~al.},
Glass-like dynamics of collective cell migration,
\href{https://doi.org/10.1073/pnas.1010059108}{Proc.\ Natl.\ Acad.\ Sci.\
USA \textbf{108}, 4714 (2011)}.

\bibitem{Garcia2015}
S.~Garcia, E.~Hannezo, J.~Elgeti, J.-F. Joanny, P.~Silberzan, and
N.~S. Gov,
Physics of active jamming during collective cellular motion in a
monolayer,
\href{https://doi.org/10.1073/pnas.1510973112}{Proc.\ Natl.\ Acad.\ Sci.\
USA \textbf{112}, 15314 (2015)}.

\bibitem{Park2015}
J.~A. Park \textit{et~al.},
Unjamming and cell shape in the asthmatic airway epithelium,
\href{https://doi.org/10.1038/nmat4357}{Nat.\ Mater.\ \textbf{14}, 1040
(2015)}.

\bibitem{Mongera2018}
A.~Mongera \textit{et~al.},
A fluid-to-solid jamming transition underlies vertebrate body axis
elongation,
\href{https://doi.org/10.1038/s41586-018-0479-2}{Nature \textbf{561}, 401
(2018)}.

\bibitem{Atia2018}
L.~Atia \textit{et~al.},
Geometric constraints during epithelial jamming,
\href{https://doi.org/10.1038/s41567-018-0089-9}{Nat.\ Phys.\ \textbf{14},
613 (2018)}.

\bibitem{Malinverno2017}
C.~Malinverno \textit{et~al.},
Endocytic reawakening of motility in jammed epithelia,
\href{https://doi.org/10.1038/nmat4848}{Nat.\ Mater.\ \textbf{16}, 587
(2017)}.

\bibitem{NagaiHonda2001}
T.~Nagai and H.~Honda,
A dynamic cell model for the formation of epithelial tissues,
\href{https://doi.org/10.1080/13642810108205772}{Philos.\ Mag.\ B
\textbf{81}, 699 (2001)}.

\bibitem{Farhadifar2007}
R.~Farhadifar, J.-C. R\"oben, B.~Aigouy, S.~Eaton, and F.~J\"ulicher,
The influence of cell mechanics, cell-cell interactions, and
proliferation on epithelial packing,
\href{https://doi.org/10.1016/j.cub.2007.11.049}{Curr.\ Biol.\
\textbf{17}, 2095 (2007)}.

\bibitem{Staple2010}
D.~B. Staple, R.~Farhadifar, J.-C. R\"oben, B.~Aigouy, S.~Eaton, and
F.~J\"ulicher,
Mechanics and remodelling of cell packings in epithelia,
\href{https://doi.org/10.1140/epje/i2010-10677-0}{Eur.\ Phys.\ J.\ E
\textbf{33}, 117 (2010)}.

\bibitem{Fletcher2014}
A.~G. Fletcher, M.~Osterfield, R.~E. Baker, and S.~Y. Shvartsman,
Vertex models of epithelial morphogenesis,
\href{https://doi.org/10.1016/j.bpj.2013.11.4498}{Biophys.\ J.\
\textbf{106}, 2291 (2014)}.

\bibitem{Bi2015}
D.~Bi, J.~H. Lopez, J.~M. Schwarz, and M.~L. Manning,
A density-independent glass transition in biological tissues,
\href{https://doi.org/10.1038/nphys3471}{Nat.\ Phys.\ \textbf{11}, 1074
(2015)}.

\bibitem{Bi2016}
D.~Bi, X.~Yang, M.~C. Marchetti, and M.~L. Manning,
Motility-driven glass and jamming transitions in biological tissues,
\href{https://doi.org/10.1103/PhysRevX.6.021011}{Phys.\ Rev.\ X
\textbf{6}, 021011 (2016)}.

\bibitem{Barton2017}
D.~L. Barton, S.~Henkes, C.~J. Weijer, and R.~Sknepnek,
Active Vertex Model for cell-resolution description of epithelial
tissue mechanics,
\href{https://doi.org/10.1371/journal.pcbi.1005569}{PLoS Comput.\ Biol.\
\textbf{13}, e1005569 (2017)}.

\bibitem{YanBi2019}
L.~Yan and D.~Bi,
Multicellular rosettes drive fluid-solid transition in epithelial
tissues,
\href{https://doi.org/10.1103/PhysRevX.9.011029}{Phys.\ Rev.\ X
\textbf{9}, 011029 (2019)}.

\bibitem{MerkelManning2018}
M.~Merkel and M.~L. Manning,
A geometrically controlled rigidity transition in a model for confluent
3D tissues,
\href{https://doi.org/10.1088/1367-2630/aab4e8}{New J.\ Phys.\
\textbf{20}, 022002 (2018)}.

\bibitem{Czajkowski2019}
M.~Czajkowski, D.~M. Sussman, M.~C. Marchetti, and M.~L. Manning,
Glassy dynamics in models of confluent tissue with mitosis and
apoptosis,
\href{https://doi.org/10.1039/C9SM00916G}{Soft Matter \textbf{15}, 9133
(2019)}.

\bibitem{Camley2017}
B.~A. Camley and W.-J. Rappel,
Physical models of collective cell motility: from cell to tissue,
\href{https://doi.org/10.1088/1361-6463/aa56fe}{J.\ Phys.\ D \textbf{50},
113002 (2017)}.

\bibitem{AlertTrepat2020}
R.~Alert and X.~Trepat,
Physical models of collective cell migration,
\href{https://doi.org/10.1146/annurev-conmatphys-031218-013516}{Annu.\
Rev.\ Condens.\ Matter Phys.\ \textbf{11}, 77 (2020)}.

\bibitem{Shao2010}
D.~Shao, W.-J. Rappel, and H.~Levine,
Computational model for cell morphodynamics,
\href{https://doi.org/10.1103/PhysRevLett.105.108104}{Phys.\ Rev.\ Lett.\
\textbf{105}, 108104 (2010)}.

\bibitem{Ziebert2012}
F.~Ziebert, S.~Swaminathan, and I.~S. Aranson,
Model for self-polarization and motility of keratocyte fragments,
\href{https://doi.org/10.1098/rsif.2011.0694}{J.\ R.\ Soc.\ Interface
\textbf{9}, 1084 (2012)}.

\bibitem{Nonomura2012}
M.~Nonomura,
Study on multicellular systems using a phase field model,
\href{https://doi.org/10.1371/journal.pone.0033501}{PLoS ONE \textbf{7},
e33501 (2012)}.

\bibitem{Lober2015}
J.~L\"ober, F.~Ziebert, and I.~S. Aranson,
Collisions of deformable cells lead to collective migration,
\href{https://doi.org/10.1038/srep09172}{Sci.\ Rep.\ \textbf{5}, 9172
(2015)}.

\bibitem{Najem2016}
S.~Najem and M.~Grant,
Phase-field model for collective cell migration,
\href{https://doi.org/10.1103/PhysRevE.93.052405}{Phys.\ Rev.\ E
\textbf{93}, 052405 (2016)}.

\bibitem{Palmieri2015}
B.~Palmieri, Y.~Bresler, D.~Wirtz, and M.~Grant,
Multiple scale model for cell migration in monolayers,
\href{https://doi.org/10.1038/srep11745}{Sci.\ Rep.\ \textbf{5}, 11745
(2015)}.

\bibitem{Loewe2020}
B.~Loewe, M.~Chiang, D.~Marenduzzo, and M.~C. Marchetti,
Solid-liquid transition of deformable and overlapping active particles,
\href{https://doi.org/10.1103/PhysRevLett.125.038003}{Phys.\ Rev.\ Lett.\
\textbf{125}, 038003 (2020)}.

\bibitem{Wenzel2021}
D.~Wenzel and A.~Voigt,
Multiphase field models for collective cell migration,
\href{https://doi.org/10.1103/PhysRevE.104.054410}{Phys.\ Rev.\ E
\textbf{104}, 054410 (2021)}.

\bibitem{Graham2024}
R.~Graham, D.~Zhang, and J.~M. Yeomans,
Cell sorting by active forces in a phase-field model of cell
monolayers,
\href{https://doi.org/10.1039/d3sm01033c}{Soft Matter \textbf{20}, 2955
(2024)}.

\bibitem{Saito2024}
S.~Saito and S.~Ishihara,
Cell deformability gives rise to a fluid-to-fluid transition in tissues,
\href{https://doi.org/10.1126/sciadv.adi8433}{Sci.\ Adv.\ \textbf{10},
eadi8433 (2024)}.

\bibitem{Moshe2018}
M.~Moshe, M.~J. Bowick, and M.~C. Marchetti,
Geometric frustration and solid-solid transitions in model 2D tissue,
\href{https://doi.org/10.1103/PhysRevLett.120.268105}{Phys.\ Rev.\ Lett.\
\textbf{120}, 268105 (2018)}.

\bibitem{Bresler2018}
Y.~Bresler, B.~Palmieri, and M.~Grant,
Sharp interface limit of the phase field model for cell motility
(2018), \href{https://arxiv.org/abs/1807.10318}{arXiv:1807.10318}.

\bibitem{HohenbergHalperin1977}
P.~C. Hohenberg and B.~I. Halperin,
Theory of dynamic critical phenomena,
\href{https://doi.org/10.1103/RevModPhys.49.435}{Rev.\ Mod.\ Phys.\
\textbf{49}, 435 (1977)}.

\end{thebibliography}

\end{document}
